\documentclass[10pt,a4paper]{amsart}
\usepackage[margin=19mm]{geometry}
\usepackage{amsmath,amssymb,mathtools}
\usepackage[scr=rsfs,cal=euler]{mathalfa}
\usepackage{xfrac}
\usepackage[unicode,hidelinks,bookmarks=false]{hyperref}
\newcommand{\ie}{i.e.\ }
\newcommand{\cf}{cf.\ }
\newcommand{\resp}{respectively\ }
\newcommand{\Subsection}{\subsection}
\providecommand{\nobreakdash}{\nobreak\hbox{-}}
\setlength{\emergencystretch}{2em}
\multlinegap0pt
\usepackage{amssymb}
\usepackage{mathtools}
\usepackage[shortlabels]{enumitem}
\usepackage{xcolor}
\newtheorem{lemma}{Lemma}
\title{A Two-Variable Zeta Function for a\\ Parity-Perturbed Hofstadter \texorpdfstring{$Q$}{Q}-Recursion\\[1.5mm] \large The Exceptional \texorpdfstring{$t=-1$}{t=-1} Slice and Gaussian Boundary Layers}
\author{Marco Mantovanelli}
\date{Source: arXiv:2609.02412v1 (CC BY 4.0)}
\begin{document}
\maketitle

\noindent\emph{Source and licence.} A Two-Variable Zeta Function for a\\ Parity-Perturbed Hofstadter \texorpdfstring{$Q$}{Q}-Recursion\\[1.5mm] \large The Exceptional \texorpdfstring{$t=-1$}{t=-1} Slice and Gaussian Boundary Layers. Version arXiv:2609.02412v1 (CC BY 4.0), distributed by arXiv under the Creative Commons Attribution 4.0 International licence, http://creativecommons.org/licenses/by/4.0/, as recorded in the arXiv licence metadata for this exact version and shown on the abstract page https://arxiv.org/abs/2609.02412v1. Full licence notice: \texttt{licenses/arxiv-2609.02412-CC-BY-4.0.txt}.\par\smallskip

\noindent\textit{Editorial scope.} Counted unit: the article's lemma lem:second-corner-coefficients (source0.tex lines 2868--2899). The article's complete original proof of it is delivered with it (source0.tex lines 2901--3336, one \texttt{proof} environment of 436 lines). No mathematics is written, repaired or re-derived: the statement and the proof are the article's own lines, in the article's own notation, and no retained line is substituted or re-flowed.  Retained uncounted as the source-local context the proof needs: lines 2517--2523; lines 2597--2605.  Nothing was omitted from any retained span, so the package carries no omission record.  No retained line contains a citation, so the wrapper carries no bibliography and no bibliography entry.  No retained line contains a figure environment, an image-inclusion command or drawing code, so no image is delivered and the package declares no asset dependency.  Editorial formatting only, in addition: an AMS \texttt{amsart} preamble that restates the article's own declarations and macros used by the retained lines, namely the environments \texttt{lemma} on separate counters, together with the standard packages those lines need.  The retained environments print in this document as Lemma 1 in order of appearance, whereas the article prints the same items with the numbering of its own full document.  The complete native source of the article as distributed in the arXiv e-print for this exact version is bundled unchanged as \texttt{sources/209196/source0.tex}.
\begin{equation}
 e_L
 =\frac12\sum_{0\le a<b<L}
 (-1)^a(b-a)
 \frac{(a+b+1)!}{(a+1)!(b+1)!}.
 \label{eq:eL-exact}
\end{equation}

\begin{align}
 \mathfrak B_r^{\rm int}
 &=\sum_{i=0}^{R}\beta_{2i},
 \label{eq:B-int-exact}\\
 \mathfrak B_r^{\rm cross}
 &=\sum_{0\le i<j\le R}(j-i)
 \frac{(2i+2j+1)!}{(2i+1)!(2j+1)!}.
 \label{eq:B-cross-exact}
\end{align}

\begin{lemma}[Second corner coefficients]
\label{lem:second-corner-coefficients}
With $R=r+1$ and $\mathscr C_R=\binom{4R}{2R}$, one has
\begin{align}
 Q_r
 &=\mathscr C_R\left(
 \frac{16}{9}-\frac{40}{81R}+\frac{61}{243R^2}
 +O(R^{-3})\right),
 \label{eq:Q-second-corner}\\
 P_r
 &=\mathscr C_R\left(
 \frac{16}{9}-\frac{112}{81R}+\frac{349}{243R^2}
 +O(R^{-3})\right),
 \label{eq:P-second-corner}\\
 \mathfrak B_r^{\rm int}
 &=\mathscr C_R\left(
 \frac{8}{405R}+\frac{2}{6075R^2}+O(R^{-3})\right),
 \label{eq:Bint-second-corner}\\
 \mathfrak B_r^{\rm cross}
 &=\mathscr C_R\left(
 \frac{64}{135R}-\frac{1184}{2025R^2}+O(R^{-3})\right).
 \label{eq:Bcross-second-corner}
\end{align}
Consequently
\begin{equation}
 {
 \mathfrak B_r
 =\mathscr C_R\left(
 \frac{40}{81R}-\frac{142}{243R^2}+O(R^{-3})\right).}
 \label{eq:B-second-corner}
\end{equation}
\end{lemma}

\begin{proof}
Put \(s=p+q\), \(d=p-q\), and let
\(\mathscr C_R=\binom{4R}{2R}\).  The exact product representation for the
first corner ratio is
\begin{equation}
 \frac{\binom{4R-2s}{2R-2p}}{\mathscr C_R}
 =4^{-s}
 \frac{
   \displaystyle\prod_{a=0}^{2p-1}\left(1-\frac{a}{2R}\right)
   \displaystyle\prod_{b=0}^{2q-1}\left(1-\frac{b}{2R}\right)}
  {\displaystyle\prod_{c=0}^{2s-1}\left(1-\frac{c}{4R}\right)}.
 \label{eq:Q-corner-product}
\end{equation}
Writing
\(S_2(m)=\sum_{j=0}^{m-1}j^2=m(m-1)(2m-1)/6\), logarithmic expansion of
\eqref{eq:Q-corner-product} gives
\begin{equation}
 \log\!\left(
 4^s\frac{\binom{4R-2s}{2R-2p}}{\mathscr C_R}\right)
 =\frac{\alpha_{p,q}}R+\frac{\lambda_{p,q}}{R^2}
 +O\!\left(\frac{(1+s)^4}{R^3}\right),
 \label{eq:Q-corner-log-second}
\end{equation}
where
\begin{equation}
 \alpha_{p,q}=-\frac{d^2}{2}+\frac{s}{4},
 \qquad
 \lambda_{p,q}
 =-\frac{S_2(2p)+S_2(2q)}8+\frac{S_2(2s)}{32}
 =-\frac{(2s-1)(4d^2-s)}{32}.
 \label{eq:alpha-lambda-corner}
\end{equation}
Exponentiating yields
\begin{equation}
 \frac{\binom{4R-2s}{2R-2p}}{\mathscr C_R}
 =4^{-s}\left(
 1+\frac{\alpha_{p,q}}R
 +\frac{\alpha^{(2)}_{p,q}}{R^2}
 +O\!\left(\frac{(1+s)^6}{R^3}\right)\right),
 \label{eq:Q-corner-second-local}
\end{equation}
with the explicit polynomial
\begin{equation}
 {
 \alpha^{(2)}_{p,q}
 =\frac{d^4}{8}-\frac{3sd^2}{8}
 +\frac{d^2}{8}+\frac{3s^2}{32}-\frac{s}{32}.}
 \label{eq:alpha-second-polynomial}
\end{equation}

The second corner ratio differs from the first by the exact factor
\begin{equation}
 \frac{\binom{4R-2s}{2R-2p+1}}
      {\binom{4R-2s}{2R-2p}}
 =\frac{2R-2q}{2R-2p+1}
 =1+\frac{d-1/2}{R}
 +\frac{(1/2-p)(q+1/2-p)}{R^2}
 +O\!\left(\frac{(1+s)^3}{R^3}\right).
 \label{eq:P-over-Q-local-factor}
\end{equation}
Consequently
\begin{equation}
 \frac{\binom{4R-2s}{2R-2p+1}}{\mathscr C_R}
 =4^{-s}\left(
 1+\frac{\beta_{p,q}}R
 +\frac{\beta^{(2)}_{p,q}}{R^2}
 +O\!\left(\frac{(1+s)^6}{R^3}\right)\right),
 \label{eq:P-corner-second-local}
\end{equation}
where
\begin{equation}
 \beta_{p,q}=\alpha_{p,q}+d-\frac12,
 \label{eq:beta-first-polynomial}
\end{equation}
and
\begin{equation}
 {
 \beta^{(2)}_{p,q}
 =\alpha^{(2)}_{p,q}
 +\alpha_{p,q}\left(d-\frac12\right)
 +\left(\frac12-p\right)\left(q+\frac12-p\right).}
 \label{eq:beta-second-polynomial}
\end{equation}

We now evaluate the two polynomially weighted geometric sums explicitly.
Put
\begin{equation}
 \mu_j:=\sum_{n\ge0}\frac{n^j}{4^n}
 =\left.
 \left(x\frac{d}{dx}\right)^j\frac1{1-x}
 \right|_{x=1/4}.
 \label{eq:quarter-geometric-moments}
\end{equation}
The required values are
\begin{equation}
 \mu_0=\frac43,\quad \mu_1=\frac49,\quad
 \mu_2=\frac{20}{27},\quad \mu_3=\frac{44}{27},\quad
 \mu_4=\frac{380}{81}.
 \label{eq:quarter-moment-values}
\end{equation}
Expanding \eqref{eq:alpha-second-polynomial} in \(p,q\) and collecting
monomials gives
\begin{align}
 \sum_{p,q\ge0}4^{-p-q}\alpha^{(2)}_{p,q}
 ={}&
 \frac14\mu_4\mu_0-\mu_3\mu_1-\frac34\mu_3\mu_0
 +\frac34\mu_2^2+\frac34\mu_2\mu_1\nonumber\\
 &+\frac7{16}\mu_2\mu_0-\frac1{16}\mu_1^2
 -\frac1{16}\mu_1\mu_0
 =\frac{61}{243}.
 \label{eq:alpha-second-geometric-sum}
\end{align}
Likewise \eqref{eq:beta-second-polynomial} gives
\begin{align}
 \sum_{p,q\ge0}4^{-p-q}\beta^{(2)}_{p,q}
 ={}&
 \frac14\mu_4\mu_0-\mu_3\mu_1-\frac34\mu_3\mu_0
 +\frac34\mu_2^2+\frac34\mu_2\mu_1\nonumber\\
 &+\frac{31}{16}\mu_2\mu_0-\frac{25}{16}\mu_1^2
 -\frac{13}{16}\mu_1\mu_0+\frac14\mu_0^2
 =\frac{349}{243}.
 \label{eq:beta-second-geometric-sum}
\end{align}
Equations \eqref{eq:alpha-second-geometric-sum} and
\eqref{eq:beta-second-geometric-sum} prove
\eqref{eq:Q-second-corner}--\eqref{eq:P-second-corner}.

\medskip
\noindent\emph{Internal-centroid contribution.}
We now supply the corner calculation underlying
\eqref{eq:Bint-second-corner}.  Write
\(\mathscr D_L:=\binom{2L}{L}\).  In the exact sum
\eqref{eq:eL-exact}, make the change of variables
\[
 a=L-1-p,\qquad b=L-1-q.
\]
Thus \(0\le q<p\le L-1\), \(d:=p-q\ge1\), \(s:=p+q\), and
\[
 (-1)^a=(-1)^{L-1}(-1)^p,\qquad b-a=d.
\]
The factorial quotient has the exact product form
\begin{align}
 \frac{(a+b+1)!}{(a+1)!(b+1)!\mathscr D_L}
 &=\frac{(2L-1-s)!}{(L-p)!(L-q)!\mathscr D_L}\nonumber\\
 &=\frac{2^{-s-1}}{L}
 \frac{\displaystyle
   \prod_{u=0}^{p-1}\left(1-\frac{u}{L}\right)
   \prod_{v=0}^{q-1}\left(1-\frac{v}{L}\right)}
 {\displaystyle
   \prod_{h=0}^{s}\left(1-\frac{h}{2L}\right)}.
 \label{eq:eL-corner-product}
\end{align}
Logarithmic expansion gives
\begin{equation}
 \frac{(a+b+1)!}{(a+1)!(b+1)!\mathscr D_L}
 =\frac{2^{-s-1}}{L}
 \left(
 1+\frac{A_{p,q}}{L}+\frac{B_{p,q}}{L^2}
 +O\!\left(\frac{(1+s)^6}{L^3}\right)\right),
 \label{eq:eL-corner-expansion}
\end{equation}
uniformly, for example, on \(s\le L^{1/8}\).  The two corner
polynomials are
\begin{align}
 A_{p,q}
 &=-\sum_{u=0}^{p-1}u-\sum_{v=0}^{q-1}v
 +\frac12\sum_{h=0}^{s}h
 =\frac{3s-d^2}{4},
 \label{eq:eL-first-corner-polynomial}\\
 B_{p,q}
 &=\frac12A_{p,q}^2
 -\frac12\sum_{u=0}^{p-1}u^2
 -\frac12\sum_{v=0}^{q-1}v^2
 +\frac18\sum_{h=0}^{s}h^2\nonumber\\
 &=\frac{d^4-10d^2s+4d^2+15s^2-2s}{32}.
 \label{eq:eL-second-corner-polynomial}
\end{align}
Since \(p=q+d\), these become
\begin{align}
 A_{q+d,q}
 &=-\frac{d^2}{4}+\frac{3d}{4}+\frac{3q}{2},
 \label{eq:eL-A-qd}\\
 B_{q+d,q}
 &=\frac{d^4}{32}-\frac{5d^3}{16}
 -\frac{5d^2q}{8}+\frac{19d^2}{32}
 +\frac{15dq}{8}-\frac{d}{16}
 +\frac{15q^2}{8}-\frac{q}{8}.
 \label{eq:eL-B-qd}
\end{align}

Substitution into \eqref{eq:eL-exact} gives
\begin{equation}
 e_L=(-1)^L\frac{\mathscr D_L}{L}
 \left(E_0+\frac{E_1}{L}+\frac{E_2}{L^2}+O(L^{-3})\right),
 \label{eq:eL-three-coefficients-pre}
\end{equation}
where, for \(F=1,A,B\), the relevant linear functional is
\begin{equation}
 \mathcal T[F]
 :=-\frac14\sum_{q\ge0}\sum_{d\ge1}
 d\left(-\frac14\right)^q
 \left(-\frac12\right)^dF(q+d,q),
 \qquad
 (E_0,E_1,E_2)=(\mathcal T[1],\mathcal T[A],\mathcal T[B]).
 \label{eq:eL-corner-functional}
\end{equation}
Put
\[
 \mathsf Q_m:=\sum_{q\ge0}q^m\left(-\frac14\right)^q,
 \qquad
 \mathsf D_m:=\sum_{d\ge1}d^m\left(-\frac12\right)^d.
\]
The needed moments are
\begin{align*}
 \mathsf Q_0&=\frac45,&
 \mathsf Q_1&=-\frac4{25},&
 \mathsf Q_2&=-\frac{12}{125},\\
 \mathsf D_1&=-\frac29,&
 \mathsf D_2&=-\frac2{27},&
 \mathsf D_3&=\frac2{27},&
 \mathsf D_4&=\frac{10}{81},&
 \mathsf D_5&=-\frac{14}{243}.
\end{align*}
Consequently
\[
 E_0=-\frac14\mathsf Q_0\mathsf D_1=\frac2{45},
\]
and, since
\(dA_{q+d,q}=-d^3/4+3d^2/4+3dq/2\),
\begin{equation}
 E_1=-\frac14\left(
 -\frac14\mathsf Q_0\mathsf D_3
 +\frac34\mathsf Q_0\mathsf D_2
 +\frac32\mathsf Q_1\mathsf D_1\right)
 =\frac1{675}.
 \label{eq:eL-E1-evaluation}
\end{equation}
Likewise,
\begin{align*}
 dB_{q+d,q}={}&
 \frac{d^5}{32}-\frac{5d^4}{16}
 -\frac{5d^3q}{8}+\frac{19d^3}{32}
 +\frac{15d^2q}{8}-\frac{d^2}{16}\\
 &+\frac{15dq^2}{8}-\frac{dq}{8},
\end{align*}
so
\begin{align}
 E_2=-\frac14\bigg(&
 \frac1{32}\mathsf Q_0\mathsf D_5
 -\frac5{16}\mathsf Q_0\mathsf D_4
 -\frac58\mathsf Q_1\mathsf D_3
 +\frac{19}{32}\mathsf Q_0\mathsf D_3\nonumber\\
 &+\frac{15}{8}\mathsf Q_1\mathsf D_2
 -\frac1{16}\mathsf Q_0\mathsf D_2
 +\frac{15}{8}\mathsf Q_2\mathsf D_1
 -\frac18\mathsf Q_1\mathsf D_1\bigg)
 =-\frac{109}{6075}.
 \label{eq:eL-E2-evaluation}
\end{align}
Thus
\begin{equation}
 {
 e_L=(-1)^L\frac1L\binom{2L}{L}
 \left(
 \frac2{45}+\frac{1}{675L}
 -\frac{109}{6075L^2}+O(L^{-3})\right).}
 \label{eq:eL-second}
\end{equation}

We next solve the exact renewal relation
\[
 \beta_L=e_L+\sum_{j=1}^{L-1}2^{j-1}e_{L-j}.
\]
Set \(G_L:=(-1)^LL^{-1}\binom{2L}{L}\).  For fixed \(j\), a
factorial-ratio expansion gives
\begin{equation}
 \frac{G_{L-j}}{G_L}
 =(-1)^j4^{-j}
 \left(1+\frac{3j}{2L}
 +\frac{j(15j-1)}{8L^2}
 +O\!\left(\frac{(1+j)^6}{L^3}\right)\right).
 \label{eq:renewal-central-ratio}
\end{equation}
Indeed,
\[
 \frac{\binom{2L-2j}{L-j}}{\binom{2L}{L}}
 =4^{-j}\left(1+\frac{j}{2L}
 +\frac{j(3j-1)}{8L^2}
 +O\!\left(\frac{(1+j)^6}{L^3}\right)\right),
\]
and multiplication by \(L/(L-j)\) gives
\eqref{eq:renewal-central-ratio}.

Write
\[
 c_0=\frac2{45},\qquad c_1=\frac1{675},\qquad
 c_2=-\frac{109}{6075},
\]
and put \(\rho_0:=1\),
\(\rho_j:=(-1)^j2^{-j-1}\) for \(j\ge1\).  Then
\begin{equation}
 W_0:=\sum_{j\ge0}\rho_j=\frac56,
 \qquad
 W_1:=\sum_{j\ge1}j\rho_j=-\frac19,
 \qquad
 W_2:=\sum_{j\ge1}j^2\rho_j=-\frac1{27}.
 \label{eq:renewal-weight-moments}
\end{equation}
Using \((L-j)^{-1}=L^{-1}+jL^{-2}+O((1+j)^2L^{-3})\)
together with \eqref{eq:renewal-central-ratio} gives
\[
 \beta_L=G_L\left(b_0+\frac{b_1}{L}+\frac{b_2}{L^2}
 +O(L^{-3})\right),
\]
where
\begin{align*}
 b_0&=c_0W_0=\frac1{27},\\
 b_1&=c_1W_0+\frac32c_0W_1=-\frac1{162},\\
 b_2&=c_2W_0+\frac52c_1W_1
 +\frac{c_0}{8}(15W_2-W_1)=-\frac{13}{729}.
\end{align*}
Therefore
\begin{equation}
 {
 \beta_L=(-1)^L\frac1L\binom{2L}{L}
 \left(
 \frac1{27}-\frac1{162L}-\frac{13}{729L^2}
 +O(L^{-3})\right).}
 \label{eq:betaL-second}
\end{equation}

Finally, write \(i=R-p\) in
\(\mathfrak B_r^{\rm int}=\sum_{i=0}^{R}\beta_{2i}\).  Uniformly on a
fixed corner window,
\[
 \frac{\binom{4(R-p)}{2(R-p)}}{\mathscr C_R}
 =16^{-p}\left(1+\frac{p}{2R}
 +O\!\left(\frac{(1+p)^2}{R^2}\right)\right).
\]
With \(b_0=1/27\), \(b_1=-1/162\), this yields
\[
 \frac{\beta_{2(R-p)}}{\mathscr C_R}
 =16^{-p}\left[
 \frac{b_0}{2R}+\frac{3pb_0+b_1}{4R^2}
 +O\!\left(\frac{(1+p)^2}{R^3}\right)\right].
\]
Since
\[
 \sum_{p\ge0}16^{-p}=\frac{16}{15},
 \qquad
 \sum_{p\ge0}p16^{-p}=\frac{16}{225},
\]
we obtain
\begin{align*}
 \frac{\mathfrak B_r^{\rm int}}{\mathscr C_R}
 &=\frac{b_0}{2R}\frac{16}{15}
 +\frac1{4R^2}\left(
 3b_0\frac{16}{225}+b_1\frac{16}{15}\right)
 +O(R^{-3})\\
 &=\frac{8}{405R}+\frac{2}{6075R^2}+O(R^{-3}),
\end{align*}
which is \eqref{eq:Bint-second-corner}.

All termwise sums in this internal-centroid calculation are uniform.  In
\eqref{eq:eL-corner-product}, Taylor expansion on
\(s\le L^{1/8}\) has the summable error displayed in
\eqref{eq:eL-corner-expansion}.  On the complement, the standard uniform
binomial estimate and \(\mathscr D_L\asymp4^L/\sqrt L\) reduce the tail,
relative to \(\mathscr D_L/L\), to a polynomially weighted geometric
tail \(\sum_{s>L^{1/8}}s^2 2^{-s}\).  Similarly,
\(|e_n|\ll\binom{2n}{n}/n\), so the range \(j>L^{1/8}\) in the renewal
sum is geometrically small relative to \(G_L\).  Finally,
\(|\beta_{2i}|\ll\binom{4i}{2i}/i\), which gives the same conclusion for
the last \(p\)-corner.  This justifies every expansion through the stated
order.

For the cross term, the factorial ratio in \eqref{eq:B-cross-exact}, with
\(i=R-p\), \(j=R-q\), factors as
\begin{align}
 &\frac{(4R-2p-2q+1)!}
 {(2R-2p+1)!(2R-2q+1)!\mathscr C_R}\nonumber\\
 &\qquad=
 \frac{4^{-p-q}}R\left(
 1+\frac{\gamma_{p,q}}R
 +O\!\left(\frac{(1+p+q)^4}{R^2}\right)\right),
 \label{eq:cross-corner-local-second}
\end{align}
where
\begin{equation}
 {
 \gamma_{p,q}
 =-\frac{p^2+q^2}{2}+pq+\frac{3(p+q)}4-\frac34.}
 \label{eq:gamma-cross-polynomial}
\end{equation}
Indeed, divide the first corner ratio
\eqref{eq:Q-corner-second-local} by the two missing linear factors; the
additional relative coefficient is \((p+q)/2-3/4\).

To sum \eqref{eq:gamma-cross-polynomial}, write \(p=q+d\), \(d\ge1\),
and set \(x=1/4\).  Then
\[
 \gamma_{q+d,q}
 =-\frac{d^2}{2}+\frac{3q}{2}+\frac{3d}{4}-\frac34.
\]
Using \eqref{eq:quarter-geometric-moments} and elementary geometric
summation gives
\begin{equation}
 \sum_{p>q\ge0}(p-q)4^{-p-q}
 =\frac{\mu_1}{1-x^2}=\frac{64}{135}.
 \label{eq:cross-leading-geometric-sum}
\end{equation}
\begin{align}
 \sum_{p>q\ge0}(p-q)4^{-p-q}\gamma_{p,q}
 &=\frac1{1-x^2}
 \left(-\frac12\mu_3+\frac34\mu_2-\frac34\mu_1\right)
 +\frac32\frac{x^2}{(1-x^2)^2}\mu_1\nonumber\\
 &=-\frac{1184}{2025}.
 \label{eq:cross-second-geometric-sum}
\end{align}
This proves \eqref{eq:Bcross-second-corner}; addition gives
\eqref{eq:B-second-corner}.

For completeness, all termwise corner summations above are uniform.  On
the window \(p+q\le R^{1/8}\), Taylor's formula gives the displayed
remainders.  Outside that window,
\[
 \frac{\binom{4R-2p-2q}{\,\cdot\,}}{\mathscr C_R}
 \ll \sqrt R\,4^{-p-q}
\]
by \(\binom Nm\le2^N\) and
\(\mathscr C_R\gg16^R/\sqrt R\); the cross ratio is bounded in the same
way.  Hence the complementary tails are smaller than every fixed inverse
power of \(R\), and the polynomially weighted geometric sums of the local
remainders converge.  This justifies expansion and summation through the
asserted orders.
\end{proof}

\end{document}
